\section{Introduction}
\label{sec:intro}

% Rule: every claim about prior work carries what OUR measurement does to it.
% The long form of this section is docs/lineage.md; do not restate it, cite it.

\todo[inline]{Draft. The argument: (i) the problem class; (ii) the aggregate is a selection
operator; (iii) the bias is known in every field under a different name; (iv) what is new
here is the empirical, per-entity, count-indexed form.}

\subsection{The problem class}
\label{sec:problem}

Many candidate entities. Each carries a noisy aggregate score estimated from a varying number
of observations. Confounds correlate with that score. A downstream validation is expensive and
can test only a handful.

\subsection{The same bias, once per field}
\label{sec:names}

The mechanism is textbook order statistics. It is rediscovered under a local name in each
field it damages: the winner's curse in genetics, the deflated Sharpe ratio in
finance~\cite{bailey2014deflated}, regression to the mean in clinical trials, and
scan-statistic inflation in change-point detection~\cite{darling1956limit}.

\subsection{What is new}
\label{sec:contribution}

Three things, stated so they can be disputed:
\begin{enumerate}
  \item The bias is estimated \emph{empirically}, from the screen's own controls, so it
        inherits the real correlation structure instead of assuming independence.
  \item It is indexed \emph{per entity by observation count}, which is the form that breaks
        comparability --- not per discovery, which is the usual framing.
  \item The direction of the bias with $\nobs$ is not assumed. It rises for some statistics
        and falls for others; see \cref{sec:direction}.
\end{enumerate}
